% diagrams/firstwords/bang.tex -- un cracker de Navidad que dos manos
% tiran de los extremos: ¡bang! Para "completa": se colorea y se dibuja
% lo que sale de dentro (una corona de papel, un chiste, un juguete).
\begin{tikzpicture}[dibujofw]
  % la mitad de la izquierda y la mano
  \filldraw[fill=white] (-4.6,-0.5) -- (-3.6,-0.3) -- (-3.6,0.3) -- (-4.6,0.5) -- cycle;
  \filldraw[fill=white] (-3.6,-0.3) -- (-3.2,-0.7) -- (-0.8,-0.7) -- (-0.8,0.7) -- (-3.2,0.7) -- (-3.6,0.3) -- cycle;
  \draw (-2.8,0.7) -- (-2.8,-0.7); \draw (-1.3,0.7) -- (-1.3,-0.7);
  \begin{scope}[shift={(-2.05,0)}, scale=0.4] \estrellaFW \end{scope}
  \filldraw[fill=white] (-5.1,0) circle (0.55);
  % la mitad de la derecha y la mano
  \filldraw[fill=white] (4.6,-0.5) -- (3.6,-0.3) -- (3.6,0.3) -- (4.6,0.5) -- cycle;
  \filldraw[fill=white] (3.6,-0.3) -- (3.2,-0.7) -- (0.8,-0.7) -- (0.8,0.7) -- (3.2,0.7) -- (3.6,0.3) -- cycle;
  \draw (2.8,0.7) -- (2.8,-0.7); \draw (1.3,0.7) -- (1.3,-0.7);
  \begin{scope}[shift={(2.05,0)}, scale=0.4] \estrellaFW \end{scope}
  \filldraw[fill=white] (5.1,0) circle (0.55);
  % el bang
  \filldraw[fill=white] (0,1.6) -- (0.35,0.6) -- (1.0,1.1) -- (0.6,0.25) -- (1.2,-0.3) -- (0.45,-0.35)
    -- (0.5,-1.3) -- (0,-0.55) -- (-0.5,-1.3) -- (-0.45,-0.35) -- (-1.2,-0.3) -- (-0.6,0.25)
    -- (-1.0,1.1) -- (-0.35,0.6) -- cycle;
  \foreach \a in {60,90,120,-60,-90,-120} {\draw (\a:1.9) -- (\a:2.5);}
\end{tikzpicture}
